% \input{\pSections "ssec-pendulum"}

%%    %%    %%    %%    %%    %%    %%    %%    %%    %%
\subsection{The Pendulum}

\begin{frame}\frametitle{\href{https://pubs.aip.org/physicstoday}{Physics Today}}
\center
	%\href{https://pubs.aip.org/physicstoday/article-abstract/64/9/42/413713/Dynamic-similarity-the-dimensionless}{
	\href{run:../local/articles-similarity/dynamic-similarity.pdf}{
	\begin{overpic}[ scale = 0.45 ]
		{\pLocalGraphics pt-dynamic-similarity}
		%\put(50, 50)	{\colorbox{white}{$a+b$}}
	\end{overpic}}
	\ \\[10pt]
	\cite{bolster2011dynamic}
\end{frame}

%
\begin{frame}\frametitle{\href{https://scienceworld.wolfram.com/physics/BuckinghamsPiTheorem.html}{Buckingham's Pi Theory}}
	%
\center
\href{https://en.wikipedia.org/wiki/Buckingham_\%CF\%80_theorem}{
Physical laws are \bl{independent} of the form of the \bl{units}
}
	%
\end{frame}
%
%
\begin{frame}\frametitle{Dimensional Analysis and the \yl{Pi Theorem}}
\center
	%\href{https://www-sciencedirect-com.libproxy.unm.edu/science/article/pii/0024379582902294}{
	\href{https://www.sciencedirect.com/science/article/pii/0024379582902294}{
	\begin{overpic}[ scale = 0.75 ]
		{\pLocalGraphics curtis-pi-03}
		%\put(50, 50)	{\colorbox{white}{$a+b$}}
	\end{overpic}}
	\ \\
	\cite{curtis1982dimensional}
\end{frame}
%
\begin{frame}\frametitle{Dimensional Analysis and the \yl{Pi Theorem}}
\center
	\emph{The \bl{most fundamental result} in dimensional analysis\\ is the (Buckingham) \bl{pi theorem}.}\\[20pt]
	\pause
	\emph{One of the \bl{most famous} examples of dimensional reasoning\\ was the derivation, by \bl{G. I. Taylor}, of the formula}
		%
\begin{equation}
	r= t^{2/5}\paren{\frac{E}{\rho_{0}}}^{1/5} f(\gamma)
\label{eq:curtis-example}
\end{equation}
		%
\end{frame}
%
\begin{frame}\frametitle{\pititle}
\center 
Assumption: $\exists$ physical law of the form 
\begin{equation}
	g(t, T, \rho_{0}, E, \gamma) = 0
\label{eq:curtis-Assumption}
\end{equation}
\end{frame}
%
\begin{frame}\frametitle{\pititle}
\center 
Two dimensionless quantities
\href{https://pdf.sciencedirectassets.com/271586/1-s2.0-S0024379500X02728/1-s2.0-0024379582902294/main.pdf\#page=2}{
\begin{equation}
	\Pi_{1} = \gamma, \quad \Pi_{2} = \frac{r^{5}\rho_{0}} {t^{2} E}
\label{eq:pi}
\end{equation}
}
\end{frame}
%
%
\begin{frame}\frametitle{Dimensional Analysis and the Pi Theorem}
	\begin{definition}[Physical Law]
		A physical law  $f( Q_{i}, . . . , Q_{m}) = 0$ relates dimensional quantities $Q_{i}, . . . , Q_{m}$. 
	\end{definition}
\end{frame}
%
\begin{frame}\frametitle{Dimensional Analysis and the Pi Theorem}
	\input{\pLocalTheorems "thm-pi"}
\end{frame}
%
%
\begin{frame}\frametitle{Similarity}
\begin{enumerate}
	\item Nuclear Explosions (\cite{taylor1950formationI, taylor1950formationII})
	\item Astrophysics (\cite{ostriker1988astrophysical})
	\item Laser driven (\cite{edwards2001investigation})
	\item Plasma (\cite{moore2005tailored})
	\item Granular (\cite{boudet2013unstable, boudet2009blast})
\end{enumerate}
\end{frame}
%

\begin{frame}\frametitle{\href{https://www.genealogy.math.ndsu.nodak.edu/id.php?id=102065}{Sedov}'s Formal Development }
\center
	\href{https://www.taylorfrancis.com/books/mono/10.1201/9780203739730/similarity-dimensional-methods-mechanics-sedov}{
	\begin{overpic}[ scale = 0.35 ]
		{\pLocalGraphics sedov-toc}
		%\put(50, 50)	{\colorbox{white}{$a+b$}}
	\end{overpic}}
\end{frame}

%%    %%    %%    %%    %%    %%    %%    %%    %%    %%
%\subsection{Taylor's Papers}
%
%%    %%    %%    %%    %%    %%    %%    %%    %%    %%
%\subsection{Similarity and dimensional methods in mechanics}
\begin{frame}\frametitle{Equations of Motion}
\center
	\href{https://archive.org/details/sedov-similarity-and-dimensional-methods-in-mechanics}{
	\begin{overpic}[ scale = 0.35 ]
		{\pLocalGraphics sedov-148}
		%\put(50, 50)	{\colorbox{white}{$a+b$}}
	\end{overpic}}
	\ \\
	\cite{sedov2018similarity}
\end{frame}
%
\begin{frame}\frametitle{Units: Familiar Quantities}
\begin{enumerate}
	\item velocity
	\item momentum
	\item kinetic energy
	\item potential energy (gravitational)
	\item force
	\item pressure
	\item density (volumetric)
\end{enumerate}
\end{frame}
%
\begin{frame}\frametitle{Units: Familiar Quantities}
\begin{table}[htp]
%\caption{default}
\begin{center}
	\begin{tabular}{lll}
			%
		velocity & $v$ & $L T^{-1}$ \\[5pt]
			%
		momentum & $p = mv$ & $M L T^{-1}$ \\[5pt]
			%
		kinetic energy & $K=\tfrac{1}{2} m v^{2}$ & $M L^{2} T^{-2}$ \\[5pt]
			%
		potential energy & $U = mgh$ & $M L^{2} T^{-2}$ \\[5pt]
			%
		force & $F = ma$ & $M L T^{-2}$ \\[5pt]
			%
		pressure & $P = F / A$ & $M T^{-2}$ \\[5pt]
			%
		density & $\delta$ & $M L^{3}$ \\[5pt]
			%
	\end{tabular}
\end{center}
\label{tab:units}
\end{table}%
\end{frame}

%
\begin{frame}\frametitle{Input Parameters with Dimension}
\begin{enumerate}
	\item $E$, energy
	\item $\rho_{0}$, air pressure
\end{enumerate}
\end{frame}
%
%
\begin{frame}\frametitle{\href{https://en.wikipedia.org/wiki/Rankine\%E2\%80\%93Hugoniot_conditions}{Rankine-Hugoniot} Jump Conditions}
\center
	Connects physical states across shock front
\end{frame}
%
\begin{frame}\frametitle{\href{https://en.wikipedia.org/wiki/Rankine\%E2\%80\%93Hugoniot_conditions}{Rankine-Hugoniot} Jump Conditions}
% https://physics.stackexchange.com/questions/611651/why-does-entropy-jump-across-a-shockwave
\center
	Conservation Laws:\\[20pt]
\begin{table}[htp]
%\caption{default}
\begin{center}
\begin{tabular}{lrcl}
	\bl{mass} & $\rho_{1} u_{1}$ & $=$ & $\rho_{2} u_{2}$\\[10pt]
	\bl{momentum} & $\rho_{1} u_{1}^{2} + p_{1}$ & $=$ & $\rho_{2} u_{2}^{2} + p_{2}$\\[10pt]
	\bl{energy} & $h_{1} + \tfrac{1}{2} u_{1}^{2}$ & $=$ & $h_{2} + \tfrac{1}{2} u_{2}^{2}$
\end{tabular}
\end{center}
\label{tab:conserved}
\end{table}%
\ \\
\pause
Not conserved: \rd{\href{https://physics.stackexchange.com/questions/611651/why-does-entropy-jump-across-a-shockwave}{Entropy}}
\end{frame}
%
%
\begin{frame}\frametitle{\href{https://en.wikipedia.org/wiki/Rankine\%E2\%80\%93Hugoniot_conditions}{Rankine-Hugoniot} Jump Conditions}
% https://physics.stackexchange.com/questions/611651/why-does-entropy-jump-across-a-shockwave
\center
	Conservation Laws:\\[20pt]
\begin{table}[htp]
%\caption{default}
\begin{center}
\begin{tabular}{lrclcrcl}
	mass & \mg{$\rho_{1} u_{1}$} & \mg{$=$} & \mg{$\rho_{2} u_{2}$} & $\Rightarrow$ & $u_{k}$ & $=$ & $\frac{m} {\rho_{k}}$\\[10pt]
	mom. & \mg{$\rho_{1} u_{1}^{2} + p_{1}$} & \mg{$=$} & \mg{$\rho_{2} u_{2}^{2} + p_{2}$} & $\Rightarrow$ & $p_{2} - p_{1}$ & $=$ & $m \paren{u_{1} - u_{2}}$\\[10pt]
	energy & \mg{$h_{1} + \tfrac{1}{2} u_{1}^{2}$} & \mg{$=$} & \mg{$h_{2} + \tfrac{1}{2} u_{2}^{2}$} & $\Rightarrow$  & $h_{2} - h_{1}$ & $=$ & $\tfrac{1}{2} \paren{u_{1}^{2} - u_{2}^{2}}$
\end{tabular}
\end{center}
\label{tab:conserved-units}
\end{table}%
\end{frame}
%
\begin{frame}\frametitle{\href{https://en.wikipedia.org/wiki/Rankine\%E2\%80\%93Hugoniot_conditions}{Rankine-Hugoniot} Jump Conditions}
\center
	Implications:
\begin{equation}
	-m^{2} = \frac{p_{2} - p_{1}} { 1 / \rho_{2} - 1 / \rho_{1} }
\label{eq:Michelson–Rayleigh}
\end{equation}
%
\begin{equation}
	h_{2} - h_{1} = \frac{1}{2} \displaystyle\paren{ 1 / \rho_{2} - 1 / \rho_{1} } \paren{p_{2} - p_{1}}
\label{eq:enthalpy}
\end{equation}
\end{frame}

%\section{Jump Conditions}
\begin{frame}\frametitle{Jump Conditions: Physical Properties}
\begin{table}[htp]
\caption{default}
\begin{center}
\begin{tabular}{lcll}
	$u_{l}$ && $u_{r}$ & density\\
	$\rho_{l}$ && $\rho_{r}$ & velocity\\
	$T_{l}$ && $T_{r}$ & pressure\\
\end{tabular}
\end{center}
\label{tab:}
\end{table}%
\end{frame}

%\section{Shock Hydrodynamics}
%\begin{frame}\frametitle{Shock Hydrodynamics}
%\center
%	\href{http://homepage.physics.uiowa.edu/~ghowes/teach/phys7729/docs/Taylor50a.pdf}{
%	\begin{overpic}[ scale = 0.5 ]
%		{\pLocalGraphics ss/taylor-01}
%		%\put(50, 50)	{\colorbox{white}{$a+b$}}
%	\end{overpic}}
%	\cite[LANL]{bethe1944shock}
%\end{frame}

\begin{frame}\frametitle{Physical Properties}
\begin{table}[htp]
\caption{default}
\begin{center}
\begin{tabular}{lcll}
	$u_{l}$ && $u_{r}$ & density\\
	$\rho_{l}$ && $\rho_{r}$ & velocity\\
	$T_{l}$ && $T_{r}$ & pressure\\
\end{tabular}
\end{center} 
\label{default}
\end{table}%
\end{frame}

\begin{frame}\frametitle{Assumptions, Simplifications}
\begin{enumerate}
	\item point explosion
	\item infinitesimal time
	\item high energy explosion
	\item infinitely strong shock
	\item 
\end{enumerate}
\end{frame}

\begin{frame}\frametitle{Assumptions, Simplifications}
\begin{enumerate}
	\item spherically symmetric blast wave
	\item point source
	\item infinitesimal time
	\item infinitely strong shock
	\item 
\end{enumerate}
\end{frame}

%
\begin{frame}\frametitle{Self-Similar Region}
\begin{enumerate}
	\item pressure behind shock wave dominates air pressure $\rho_{1} >> \rho_0$
	\item point source
	\item infinitesimal time
	\item infinitely strong shock
	\item 
\end{enumerate}
\end{frame}
%

\endinput  %  ==  ==  ==  ==  ==  ==  ==  ==  ==
